\chapter{The LeanFM Web-Server Protocol}
\label{ch:web-protocol}

\index{HTTP!protobuf adapter}
\index{protocol!LeanFM web server}
This chapter is the protocol reference for the LeanFM web server. It is
primarily machine-oriented context for an implementation agent and may be read
as a maintained portion of the system prompt. The normative behavioral
source is \nolinkurl{LeanFM/ChatUI/Requirements.lean}; the byte vocabulary is
\nolinkurl{LeanFM/ChatUI/Requirements.proto}; and transport selection is
\nolinkurl{LeanFM/ChatUI/Implementation.lean}. The tables and grammars below are
generated from those files. They are not a second, hand-maintained API design.

Two versions of the protocol must not be confused. The \emph{accepted protocol}
is the complete typed grammar and intended HTTP mapping. The \emph{running
adapter} is the subset currently dispatched by \texttt{LeanFM/ChatUI/Server.lean}.
A route in the accepted table is not a claim that its handler already exists.

\section{One exchange across four representations}

\begin{figure}[htbp]
\centering
\resizebox{\textwidth}{!}{% Generated by LeanFM.MDPDiagram.renderInteractionTeX.
\begin{tikzpicture}[actor/.style={draw=diagramink,text=diagramink,fill=diagramsurface,rounded corners,minimum width=2.3cm,minimum height=7mm,font=\small},message/.style={->,>=Stealth,draw=leanblue,thick}]
\node[actor] (a0) at (0,0) {Browser};
\draw[dashed,draw=diagramink] (a0.south) -- ++(0,-6);
\node[actor] (a1) at (6,0) {LeanFM Server};
\draw[dashed,draw=diagramink] (a1.south) -- ++(0,-6);
\node[actor] (a2) at (12,0) {Grammar};
\draw[dashed,draw=diagramink] (a2.south) -- ++(0,-6);
\node[actor] (a3) at (18,0) {Event Journal};
\draw[dashed,draw=diagramink] (a3.south) -- ++(0,-6);
\draw[message] (0,-1) -- node[above,fill=pagebackground,text=diagramink,font=\scriptsize] {HTTP request(method, path, fields, session cookie)} (6,-1);
\draw[message] (6,-2) -- node[above,fill=pagebackground,text=diagramink,font=\scriptsize] {construct typed request payload} (12,-2);
\draw[message] (12,-3) -- node[above,fill=pagebackground,text=diagramink,font=\scriptsize] {append EventEnvelope(payload, priorIds)} (18,-3);
\draw[message] (12,-4) -- node[above,fill=pagebackground,text=diagramink,font=\scriptsize] {typed response payload or rejection} (6,-4);
\draw[message] (6,-5) -- node[above,fill=pagebackground,text=diagramink,font=\scriptsize] {HTTP status, JSON projection, Set-Cookie if applicable} (0,-5);
\end{tikzpicture}
}
\caption{HTTP is an adapter around the typed LeanFM conversation. Source,
reducer, question, and unit are recorded in
\texttt{generated/web-protocol-provenance.txt}.}
\label{fig:web-protocol-overview}
\end{figure}

The figure summarizes the round trip: authenticate and decode HTTP into a
grammar-permitted protobuf value, record its \texttt{EventEnvelope}, then
project the selected response or rejection back to HTTP. Provider messages
remain in-process, not invented HTTP calls. Protobuf fixes the typed value and
stable tags; HTTP chooses the browser representation, so their bytes differ.
The current adapter uses \texttt{application/x-www-form-urlencoded} for POST
requests and JSON for responses; session establishment also issues a secure
cookie. The journal stores length-delimited binary \texttt{EventEnvelope}s. A
future binary HTTP codec may change representation without changing message
values or correlation fields. Field-level conversion follows.

\subsection{Conservative default HTTP projection}

LeanFM can generate an HTTP projection and its OpenAPI description when
\texttt{Implementation.lean} does not provide an explicit binding. The task
grammar supplies request/response pairing: an initiating browser-to-server
message is the request, its correlated finite server-to-browser alternatives
are responses, and their protobuf schemas supply the operation schemas. An
explicit implementation binding always overrides the default.

The default is deliberately conservative. A field may appear in a path or
query only when the implementation explicitly classifies it as a bounded,
URL-safe scalar selector. Bytes, repeated fields, nested messages, secrets,
large or unbounded strings, and any request whose arguments would be unnatural
as a URL use \texttt{POST}; the protobuf value is carried in the request body.
This prevents a convenient generator from leaking credentials, exceeding
practical URL limits, or inventing ambiguous encodings. Generated encoders and
decoders must preserve the protobuf value exactly.

Unless an explicit transport binding selects another codec, the default body
encoder and decoder use the protobuf JSON mapping with
\texttt{Content-Type: application/json}. This default is part of the generated
adapter contract, not a change to the protobuf value. The server publishes the
resolved OpenAPI document at \texttt{GET /openapi.yaml}; this is the default
discovery endpoint for generated HTTP operations and schemas.

The generated OpenAPI operation records the resolved method and path, body
content type, request schema, every grammar-permitted response schema and
status, authentication requirement, correlation fields, and LeanFM requirement
identifier. OpenAPI documents the transport projection; the protobuf and task
grammar remain the definitions of values and legal ordering.

\section{The conversion, field by field}

For a request binding \texttt{httpRequest(method, pathTemplate)}, conversion is
defined in this order:

\begin{enumerate}
  \item Match the HTTP method exactly and split the path into decoded segments.
  \item For each \texttt{\{field\}} in the template, take the protobuf field of
        that name, encode it as one path segment, and reject slashes, traversal,
        or a value that escapes the authenticated project root.
  \item For the currently implemented POST subset, encode remaining browser
        input fields as UTF-8 form values. A protobuf \texttt{bytes} field is
        opaque bytes at the protocol boundary; a form adapter must use an
        explicitly declared textual encoding and must never place credentials
        in a URL, trace, diagnostic, or returned object.
  \item Recover the authenticated user from the
        \texttt{leanfm\_session} HttpOnly, SameSite cookie. The cookie is
        transport authorization, not a field silently added to every protobuf
        message. Its resolved user and session identifiers constrain which
        project-scoped message may be constructed.
  \item Validate required fields and task \texttt{where\_} constraints. On
        success, assign the envelope metadata and the oneof payload. On failure,
        emit the named rejection alternative when the grammar has one; otherwise
        return an HTTP error without pretending a successful protocol event
        occurred.
\end{enumerate}

For a response binding \texttt{httpResponse(RequestName)}, conversion reverses
the projection. The request name selects the finite legal response set. Scalar
fields become JSON strings, booleans, or numbers; repeated fields become JSON
arrays; nested messages become objects; and \texttt{bytes} become an explicitly
identified download or Base64 field rather than accidental UTF-8. Secrets are
redacted before JSON rendering. Success commonly uses 200, creation uses 201,
authentication failure uses 401 or 403, invalid input uses 400, missing owned
objects use 404, revision or credential conflicts use 409, and provider failure
uses 502. The named response message, not the status alone, determines the
grammar transition.

Before projection, request and response are independent framed events. In that
limited structural sense the stream resembles a multiplexed protocol such as
HTTP/2: several exchanges may be in flight and correlation is explicit. This is
not a claim that the event bytes are HTTP/2 frames. Every response envelope must
set \texttt{cause\_id} to the exact request event identifier within the same
session and task. The HTTP adapter follows that causal backpointer to recover the
exchange and emit the response on the connection that carried the request. It
never pairs by adjacency, arrival order, or a process-global ``current request.''
The separate \texttt{prior\_ids} list constrains legal interleavings and is not
used as a substitute for causal correlation.

HTTP projection is valid only for a causally closed exchange. Every projected
request must be named by exactly one terminal response's \texttt{cause\_id}, and
every projected response must name one compatible request in the same session
and task. Failure is still a response event: it carries an HTTP error status and
points to the request. An orphan request, a response with no cause, an
incompatible or cross-session cause, or two terminal responses for one request
is an unchecked HTTP call and fails validation rather than becoming an
underspecified adapter behavior. Protocols that deliberately permit interim or
streaming responses must model those nonterminal messages explicitly and still
identify one terminal response.

\section{The running adapter}

The following table is the implemented ChatUI subset. Static assets,
diagnostic reports, generated source downloads, MCP, and health checks are
document resources or tool protocols outside the ChatUI task grammar; they must
not be mistaken for user-task protobuf events.

\begin{description}[style=nextline,leftmargin=0pt]
\item[POST \nolinkurl{/api/accounts}] Request: AccountCreateSubmitted. Response: AccountCreated or AccountCreateRejected. Encoding: application/x-www-form-urlencoded.
\item[POST \nolinkurl{/api/login}] Request: LoginSubmitted. Response: LoginAccepted or LoginRejected. Encoding: application/x-www-form-urlencoded.
\item[GET \nolinkurl{/api/me}] Request: authenticated session. Response: user identity JSON. Encoding: cookie only.
\item[GET \nolinkurl{/api/projects}] Request: ProjectListRequested. Response: ProjectListReturned. Encoding: cookie plus path.
\item[POST \nolinkurl{/api/projects}] Request: ProjectCreateRequested. Response: ProjectCreated. Encoding: application/x-www-form-urlencoded.
\item[GET \nolinkurl{/api/projects/{project_id}/conversation}] Request: project identifier. Response: conversation JSON. Encoding: path field plus cookie.
\item[GET \nolinkurl{/api/projects/{project_id}/tasks}] Request: project identifier. Response: task model JSON. Encoding: path field plus cookie.
\item[GET \nolinkurl{/api/projects/{project_id}/artifacts}] Request: ArtifactListRequested. Response: ArtifactListReturned. Encoding: path field plus cookie.
\item[POST \nolinkurl{/api/projects/{project_id}/chat}] Request: ChatPromptSubmitted. Response: ChatResponseRendered or ChatTurnFailed. Encoding: application/x-www-form-urlencoded.
\end{description}


Every \texttt{/api/} route except public account creation and login requires a
valid in-memory session. Project paths additionally resolve the authenticated
owner's project root. A missing session yields 401; an unknown or unauthorized
project is deliberately represented as 404. The server binds to loopback by
default. Network exposure requires an authenticating reverse proxy or an
equivalent transport boundary.

\section{The accepted HTTP projection}

This is the complete intended browser mapping from the accepted implementation
specification. Rows absent from the running-adapter table remain implementation
obligations.

\begin{description}[style=nextline,leftmargin=0pt]
\item[POST \nolinkurl{/api/login}] Request: \nolinkurl{LoginSubmitted}. Response: \nolinkurl{LoginAccepted}, \nolinkurl{LoginRejected}.
\item[POST \nolinkurl{/api/accounts}] Request: \nolinkurl{AccountCreateSubmitted}. Response: \nolinkurl{AccountCreated}, \nolinkurl{AccountCreateRejected}.
\item[GET \nolinkurl{/api/projects/{project_id}/tasks/{task_id}}] Request: \nolinkurl{TaskModelRequested}. Response: \nolinkurl{TaskModelReturned}.
\item[GET \nolinkurl{/api/projects/{project_id}/tasks/{task_id}/sessions/{session_id}/events}] Request: \nolinkurl{MessageSequenceRequested}. Response: \nolinkurl{MessageSequenceReturned}.
\item[GET \nolinkurl{/api/projects/{project_id}/events/{event_id}}] Request: \nolinkurl{ProtobufObjectRequested}. Response: \nolinkurl{ProtobufObjectReturned}.
\item[POST \nolinkurl{/api/projects/{project_id}/research}] Request: \nolinkurl{WebResearchRequested}. Response: \nolinkurl{WebResearchReturned}, \nolinkurl{WebResearchFailed}.
\item[POST \nolinkurl{/api/projects/{project_id}/markerboard/assist}] Request: \nolinkurl{MarkerboardAssistRequested}. Response: \nolinkurl{MarkerboardSyntaxHelpReturned}, \nolinkurl{MarkerboardPatchProposed}.
\item[POST \nolinkurl{/api/projects/{project_id}/markerboard/assist/{artifact_id}/accept}] Request: \nolinkurl{MarkerboardPatchAccepted}. Response: \nolinkurl{MarkerboardDraftUpdated}.
\item[POST \nolinkurl{/api/projects/{project_id}/markerboard/assist/{artifact_id}/reject}] Request: \nolinkurl{MarkerboardPatchRejected}. Response: \nolinkurl{HTTP 204 (no payload)}.
\item[GET \nolinkurl{/api/projects}] Request: \nolinkurl{ProjectListRequested}. Response: \nolinkurl{ProjectListReturned}.
\item[POST \nolinkurl{/api/projects}] Request: \nolinkurl{ProjectCreateRequested}. Response: \nolinkurl{ProjectCreated}.
\item[GET \nolinkurl{/api/projects/{project_id}}] Request: \nolinkurl{ProjectOpenRequested}. Response: \nolinkurl{ProjectOpened}, \nolinkurl{ProblemInquiryPresented}.
\item[DELETE \nolinkurl{/api/projects/{project_id}}] Request: \nolinkurl{ProjectDeleteRequested}. Response: \nolinkurl{ProjectDeleted}.
\item[POST \nolinkurl{/api/projects/{project_id}/specification/problem}] Request: \nolinkurl{ProblemStatementSubmitted}. Response: \nolinkurl{AmbiguityQuestionPresented}, \nolinkurl{RequirementsReviewPresented}.
\item[POST \nolinkurl{/api/projects/{project_id}/specification/ambiguities/{ambiguity_id}}] Request: \nolinkurl{ClarificationSubmitted}. Response: \nolinkurl{ProjectBookRevised}.
\item[POST \nolinkurl{/api/projects/{project_id}/specification/requirements/{requirements_revision}/accept}] Request: \nolinkurl{SpecificationRequirementsAccepted}. Response: \nolinkurl{ImplementationReviewPresented}.
\item[POST \nolinkurl{/api/projects/{project_id}/specification/implementation/{implementation_revision}/accept}] Request: \nolinkurl{SpecificationImplementationAccepted}. Response: \nolinkurl{ImplementationAuthorizationConfirmed}.
\item[POST \nolinkurl{/api/credentials}] Request: \nolinkurl{ProviderCredentialSubmitted}. Response: \nolinkurl{ProviderCredentialAccepted}, \nolinkurl{ProviderCredentialRejected}.
\item[POST \nolinkurl{/api/projects/{project_id}/conversations/{conversation_id}/messages}] Request: \nolinkurl{ChatPromptSubmitted}. Response: \nolinkurl{ChatResponseRendered}, \nolinkurl{ChatTurnFailed}.
\item[POST \nolinkurl{/api/projects/{project_id}/markerboard/mermaid}] Request: \nolinkurl{MermaidDraftSaved}. Response: \nolinkurl{MermaidPreviewRendered}.
\item[POST \nolinkurl{/api/projects/{project_id}/markerboard/media}] Request: \nolinkurl{MediaUploaded}. Response: \nolinkurl{MediaRendered}.
\item[POST \nolinkurl{/api/projects/{project_id}/artifacts/{artifact_id}/promotions}] Request: \nolinkurl{ArtifactPromoted}. Response: \nolinkurl{CandidateRequirementUpdated}.
\item[GET \nolinkurl{/api/projects/{project_id}/artifacts}] Request: \nolinkurl{ArtifactListRequested}. Response: \nolinkurl{ArtifactListReturned}.
\item[GET \nolinkurl{/api/projects/{project_id}/artifacts/{artifact_id}/revisions/{revision}}] Request: \nolinkurl{ArtifactOpenRequested}. Response: \nolinkurl{ArtifactOpened}.
\end{description}


Messages between \texttt{LeanFMServer} and \texttt{ModelProvider} map to the
bounded in-process channels \texttt{provider.requests} and
\texttt{provider.results}. They retain their protobuf message identities and
request identifiers even though they do not cross HTTP. This is important for
concurrency: two users can use different provider credentials because the key,
session, user, project, conversation, and request identities remain request-local
instead of being stored in a process-wide environment variable.

\section{Complete protobuf payload vocabulary}

Every accepted event is a length-delimited \texttt{EventEnvelope}. Field 20 and
above in its \texttt{payload} oneof identify the concrete message. The table
below is generated from the Lean \texttt{MessageSchema} values; the checked
\texttt{.proto} file supplies repeated and nested cardinality that the compact
Lean field summary represents only by its scalar boundary.

\begin{description}[style=nextline,leftmargin=0pt]
\item[\nolinkurl{LoginSubmitted}] UserBrowser sends to LeanFMServer.
\item[\nolinkurl{LoginAccepted}] LeanFMServer sends to UserBrowser.
\item[\nolinkurl{LoginRejected}] LeanFMServer sends to UserBrowser.
\item[\nolinkurl{AccountCreateSubmitted}] UserBrowser sends to LeanFMServer.
\item[\nolinkurl{ProviderKeyCheckRequested}] LeanFMServer sends to ModelProvider.
\item[\nolinkurl{ProviderKeyVerified}] ModelProvider sends to LeanFMServer.
\item[\nolinkurl{ProviderKeyRejected}] ModelProvider sends to LeanFMServer.
\item[\nolinkurl{AccountCreated}] LeanFMServer sends to UserBrowser.
\item[\nolinkurl{AccountCreateRejected}] LeanFMServer sends to UserBrowser.
\item[\nolinkurl{TaskModelRequested}] UserBrowser sends to LeanFMServer.
\item[\nolinkurl{TaskModelReturned}] LeanFMServer sends to UserBrowser.
\item[\nolinkurl{MessageSequenceRequested}] UserBrowser sends to LeanFMServer.
\item[\nolinkurl{MessageSequenceReturned}] LeanFMServer sends to UserBrowser.
\item[\nolinkurl{ProtobufObjectRequested}] UserBrowser sends to LeanFMServer.
\item[\nolinkurl{ProtobufObjectReturned}] LeanFMServer sends to UserBrowser.
\item[\nolinkurl{WebResearchRequested}] UserBrowser sends to LeanFMServer.
\item[\nolinkurl{WebResearchDispatched}] LeanFMServer sends to ModelProvider.
\item[\nolinkurl{WebSearchPerformed}] ModelProvider sends to LeanFMServer.
\item[\nolinkurl{WebPageOpened}] ModelProvider sends to LeanFMServer.
\item[\nolinkurl{WebResearchReturned}] LeanFMServer sends to UserBrowser.
\item[\nolinkurl{WebResearchFailed}] LeanFMServer sends to UserBrowser.
\item[\nolinkurl{MarkerboardAssistRequested}] UserBrowser sends to LeanFMServer.
\item[\nolinkurl{MarkerboardSyntaxHelpReturned}] LeanFMServer sends to UserBrowser.
\item[\nolinkurl{MarkerboardAssistDispatched}] LeanFMServer sends to ModelProvider.
\item[\nolinkurl{MarkerboardAssistReturned}] ModelProvider sends to LeanFMServer.
\item[\nolinkurl{MarkerboardPatchProposed}] LeanFMServer sends to UserBrowser.
\item[\nolinkurl{MarkerboardPatchAccepted}] UserBrowser sends to LeanFMServer.
\item[\nolinkurl{MarkerboardPatchRejected}] UserBrowser sends to LeanFMServer.
\item[\nolinkurl{MarkerboardDraftUpdated}] LeanFMServer sends to UserBrowser.
\item[\nolinkurl{ProjectListRequested}] UserBrowser sends to LeanFMServer.
\item[\nolinkurl{ProjectListReturned}] LeanFMServer sends to UserBrowser.
\item[\nolinkurl{ProjectCreateRequested}] UserBrowser sends to LeanFMServer.
\item[\nolinkurl{ProjectCreated}] LeanFMServer sends to UserBrowser.
\item[\nolinkurl{ProjectOpenRequested}] UserBrowser sends to LeanFMServer.
\item[\nolinkurl{ProjectOpened}] LeanFMServer sends to UserBrowser.
\item[\nolinkurl{ProjectDeleteRequested}] UserBrowser sends to LeanFMServer.
\item[\nolinkurl{ProjectDeleted}] LeanFMServer sends to UserBrowser.
\item[\nolinkurl{ProblemInquiryPresented}] LeanFMServer sends to UserBrowser.
\item[\nolinkurl{ProblemStatementSubmitted}] UserBrowser sends to LeanFMServer.
\item[\nolinkurl{AmbiguityQuestionPresented}] LeanFMServer sends to UserBrowser.
\item[\nolinkurl{ClarificationSubmitted}] UserBrowser sends to LeanFMServer.
\item[\nolinkurl{ProjectBookRevised}] LeanFMServer sends to UserBrowser.
\item[\nolinkurl{RequirementsReviewPresented}] LeanFMServer sends to UserBrowser.
\item[\nolinkurl{SpecificationRequirementsAccepted}] UserBrowser sends to LeanFMServer.
\item[\nolinkurl{ImplementationReviewPresented}] LeanFMServer sends to UserBrowser.
\item[\nolinkurl{SpecificationImplementationAccepted}] UserBrowser sends to LeanFMServer.
\item[\nolinkurl{ImplementationAuthorizationConfirmed}] LeanFMServer sends to UserBrowser.
\item[\nolinkurl{ProviderCredentialSubmitted}] UserBrowser sends to LeanFMServer.
\item[\nolinkurl{ProviderCredentialAccepted}] LeanFMServer sends to UserBrowser.
\item[\nolinkurl{ProviderCredentialRejected}] LeanFMServer sends to UserBrowser.
\item[\nolinkurl{ChatPromptSubmitted}] UserBrowser sends to LeanFMServer.
\item[\nolinkurl{ModelRequestIssued}] LeanFMServer sends to ModelProvider.
\item[\nolinkurl{ModelResponseReturned}] ModelProvider sends to LeanFMServer.
\item[\nolinkurl{ChatResponseRendered}] LeanFMServer sends to UserBrowser.
\item[\nolinkurl{ChatTurnFailed}] LeanFMServer sends to UserBrowser.
\item[\nolinkurl{ConversationCheckpointRequested}] UserBrowser sends to LeanFMServer.
\item[\nolinkurl{ConversationCheckpointBlocked}] LeanFMServer sends to UserBrowser.
\item[\nolinkurl{ConversationCheckpointCommitted}] LeanFMServer sends to UserBrowser.
\item[\nolinkurl{ConversationHistoryCleared}] LeanFMServer sends to UserBrowser.
\item[\nolinkurl{MermaidDraftSaved}] UserBrowser sends to LeanFMServer.
\item[\nolinkurl{MermaidPreviewRendered}] LeanFMServer sends to UserBrowser.
\item[\nolinkurl{MediaUploaded}] UserBrowser sends to LeanFMServer.
\item[\nolinkurl{MediaRendered}] LeanFMServer sends to UserBrowser.
\item[\nolinkurl{ArtifactPromoted}] UserBrowser sends to LeanFMServer.
\item[\nolinkurl{CandidateRequirementUpdated}] LeanFMServer sends to UserBrowser.
\item[\nolinkurl{ArtifactListRequested}] UserBrowser sends to LeanFMServer.
\item[\nolinkurl{ArtifactListReturned}] LeanFMServer sends to UserBrowser.
\item[\nolinkurl{ArtifactOpenRequested}] UserBrowser sends to LeanFMServer.
\item[\nolinkurl{ArtifactOpened}] LeanFMServer sends to UserBrowser.
\end{description}


\subsection{Canonical byte schema}

The following complete schema is normative. It fixes envelope and payload tag
numbers, repeated and nested cardinality, enum values, and field names. A
reconstruction must reproduce these values rather than assigning tags in book
order or replacing nested messages with opaque bytes.

\begin{lstlisting}[language={},basicstyle=\ttfamily\tiny,breaklines=true]
syntax = "proto3";

package leanfm.chatui;

message EventEnvelope {
  string event_id = 1;
  repeated string prior_ids = 2;
  string session_id = 3;
  string task_id = 4;
  uint64 time_at_ms = 5;
  string src = 6;
  string dst = 7;
  // Direct causal parent for any event, including nonterminal and intermediate
  // events. Empty for a causal root. Unlike prior_ids, this does not constrain
  // otherwise legal concurrent interleavings. Renderers may use the causal
  // subtree as nested interaction boundaries.
  optional string cause_id = 8;

  oneof payload {
    LoginSubmitted login_submitted = 20;
    LoginAccepted login_accepted = 21;
    LoginRejected login_rejected = 22;
    ProjectListRequested project_list_requested = 23;
    ProjectListReturned project_list_returned = 24;
    ProjectCreateRequested project_create_requested = 25;
    ProjectCreated project_created = 26;
    ProjectOpenRequested project_open_requested = 27;
    ProjectOpened project_opened = 28;
    ProjectDeleteRequested project_delete_requested = 29;
    ProjectDeleted project_deleted = 30;
    ProviderCredentialSubmitted provider_credential_submitted = 31;
    ProviderCredentialAccepted provider_credential_accepted = 32;
    ProviderCredentialRejected provider_credential_rejected = 33;
    ChatPromptSubmitted chat_prompt_submitted = 34;
    ModelRequestIssued model_request_issued = 35;
    ModelResponseReturned model_response_returned = 36;
    ChatResponseRendered chat_response_rendered = 37;
    ChatTurnFailed chat_turn_failed = 38;
    MermaidDraftSaved mermaid_draft_saved = 39;
    MermaidPreviewRendered mermaid_preview_rendered = 40;
    MediaUploaded media_uploaded = 41;
    MediaRendered media_rendered = 42;
    ArtifactPromoted artifact_promoted = 43;
    CandidateRequirementUpdated candidate_requirement_updated = 44;
    ArtifactListRequested artifact_list_requested = 45;
    ArtifactListReturned artifact_list_returned = 46;
    ArtifactOpenRequested artifact_open_requested = 47;
    ArtifactOpened artifact_opened = 48;
    AccountCreateSubmitted account_create_submitted = 49;
    ProviderKeyCheckRequested provider_key_check_requested = 50;
    ProviderKeyVerified provider_key_verified = 51;
    ProviderKeyRejected provider_key_rejected = 52;
    AccountCreated account_created = 53;
    AccountCreateRejected account_create_rejected = 54;
    TaskModelRequested task_model_requested = 55;
    TaskModelReturned task_model_returned = 56;
    MessageSequenceRequested message_sequence_requested = 57;
    MessageSequenceReturned message_sequence_returned = 58;
    ProtobufObjectRequested protobuf_object_requested = 59;
    ProtobufObjectReturned protobuf_object_returned = 60;
    WebResearchRequested web_research_requested = 61;
    WebSearchPerformed web_search_performed = 62;
    WebPageOpened web_page_opened = 63;
    WebResearchReturned web_research_returned = 64;
    WebResearchFailed web_research_failed = 65;
    WebResearchDispatched web_research_dispatched = 66;
    MarkerboardAssistRequested markerboard_assist_requested = 67;
    MarkerboardSyntaxHelpReturned markerboard_syntax_help_returned = 68;
    MarkerboardAssistDispatched markerboard_assist_dispatched = 69;
    MarkerboardAssistReturned markerboard_assist_returned = 70;
    MarkerboardPatchProposed markerboard_patch_proposed = 71;
    MarkerboardPatchAccepted markerboard_patch_accepted = 72;
    MarkerboardPatchRejected markerboard_patch_rejected = 73;
    MarkerboardDraftUpdated markerboard_draft_updated = 74;
    ProblemInquiryPresented problem_inquiry_presented = 75;
    ProblemStatementSubmitted problem_statement_submitted = 76;
    AmbiguityQuestionPresented ambiguity_question_presented = 77;
    ClarificationSubmitted clarification_submitted = 78;
    ProjectBookRevised project_book_revised = 79;
    RequirementsReviewPresented requirements_review_presented = 80;
    SpecificationRequirementsAccepted specification_requirements_accepted = 81;
    ImplementationReviewPresented implementation_review_presented = 82;
    SpecificationImplementationAccepted specification_implementation_accepted = 83;
    ImplementationAuthorizationConfirmed implementation_authorization_confirmed = 84;
    ConversationCheckpointRequested conversation_checkpoint_requested = 85;
    ConversationCheckpointBlocked conversation_checkpoint_blocked = 86;
    ConversationCheckpointCommitted conversation_checkpoint_committed = 87;
    ConversationHistoryCleared conversation_history_cleared = 88;
  }
}

enum ContentKind {
  CONTENT_KIND_UNSPECIFIED = 0;
  MARKDOWN = 1;
  LATEX = 2;
  MERMAID = 3;
  GRAPHVIZ_DOT = 4;
  IMAGE = 5;
  VIDEO = 6;
  CODE = 7;
  FILE = 8;
  LEANFM_VIEW = 9;
  URL_CITATION = 10;
}

message ChatContentPart {
  string part_id = 1;
  ContentKind kind = 2;
  string text = 3;
  string artifact_id = 4;
  uint64 artifact_revision = 5;
  string media_type = 6;
  bytes inline_content = 7;
  string alt_text = 8;
  string url = 9;
  string title = 10;
}

message LoginSubmitted {
  string username = 1;
  bytes credential = 2;
  bytes openai_api_key = 3;
  bytes anthropic_api_key = 4;
}
message LoginAccepted { string user_id = 1; string session_id = 2; }
message LoginRejected { string reason = 1; }

message ProblemInquiryPresented { string project_id = 1; string question = 2; }
message ProblemStatementSubmitted { string project_id = 1; string statement = 2; }
message AmbiguityQuestionPresented { string project_id = 1; string ambiguity_id = 2; string question = 3; repeated string affected_artifacts = 4; }
message ClarificationSubmitted { string project_id = 1; string ambiguity_id = 2; string answer = 3; }
message ProjectBookRevised { string project_id = 1; string book_revision = 2; repeated string decision_ids = 3; string light_pdf = 4; string dark_pdf = 5; }
message RequirementsReviewPresented { string project_id = 1; string requirements_revision = 2; repeated string derived_artifact_ids = 3; repeated string open_ambiguity_ids = 4; }
message SpecificationRequirementsAccepted { string project_id = 1; string requirements_revision = 2; }
message ImplementationReviewPresented { string project_id = 1; string requirements_revision = 2; string implementation_revision = 3; string traceability_report = 4; }
message SpecificationImplementationAccepted { string project_id = 1; string requirements_revision = 2; string implementation_revision = 3; }
message ImplementationAuthorizationConfirmed { string project_id = 1; string requirements_revision = 2; string implementation_revision = 3; }

message AccountCreateSubmitted {
  string username = 1;
  bytes login_credential = 2;
  string provider = 3;
  bytes api_key = 4;
  bytes openai_api_key = 5;
  bytes anthropic_api_key = 6;
}
message ProviderKeyCheckRequested { string provider = 1; string verification_id = 2; }
message ProviderKeyVerified { string verification_id = 1; string provider_account = 2; }
message ProviderKeyRejected { string verification_id = 1; string reason = 2; }
message AccountCreated { string user_id = 1; string credential_ref = 2; }
message AccountCreateRejected { string reason = 1; }

message TaskModelRequested { string project_id = 1; string task_id = 2; string requirements_revision = 3; }
message TaskModelReturned {
  string task_id = 1;
  string interaction_artifact_id = 2;
  string state_machine_artifact_id = 3;
  string message_sequence_artifact_id = 4;
  repeated string constraint_ids = 5;
}
message MessageSequenceRequested { string project_id = 1; string task_id = 2; string session_id = 3; }
message MessageSequenceReturned { string task_id = 1; string session_id = 2; repeated string event_ids = 3; }
message ProtobufObjectRequested { string project_id = 1; string event_id = 2; }
message ProtobufObjectReturned { string event_id = 1; string protobuf_type = 2; bytes encoded_object = 3; string redacted_json = 4; }
message WebResearchRequested { string project_id = 1; string conversation_id = 2; string request_id = 3; string question = 4; repeated string supplied_urls = 5; }
message WebResearchDispatched { string request_id = 1; string model = 2; string research_mode = 3; }
message WebSearchPerformed { string request_id = 1; repeated string queries = 2; }
message WebPageOpened { string request_id = 1; string url = 2; string title = 3; }
message WebResearchReturned { string request_id = 1; repeated ChatContentPart parts = 2; repeated string source_urls = 3; }
message WebResearchFailed { string request_id = 1; string reason = 2; }
message MarkerboardAssistRequested { string project_id = 1; string artifact_id = 2; uint64 base_revision = 3; string diagram_kind = 4; string source = 5; string instruction = 6; string mode = 7; string credential_ref = 8; }
message MarkerboardSyntaxHelpReturned { string artifact_id = 1; uint64 base_revision = 2; string help_markdown = 3; }
message MarkerboardAssistDispatched { string request_id = 1; string artifact_id = 2; uint64 base_revision = 3; string diagram_kind = 4; string source = 5; string instruction = 6; }
message MarkerboardAssistReturned { string request_id = 1; string proposed_source = 2; string explanation = 3; uint64 input_tokens = 4; uint64 output_tokens = 5; }
message MarkerboardPatchProposed { string artifact_id = 1; uint64 base_revision = 2; string proposed_source = 3; string explanation = 4; }
message MarkerboardPatchAccepted { string artifact_id = 1; uint64 base_revision = 2; string proposed_source = 3; }
message MarkerboardPatchRejected { string artifact_id = 1; uint64 base_revision = 2; }
message MarkerboardDraftUpdated { string artifact_id = 1; uint64 revision = 2; string source = 3; }

message ProjectListRequested { string user_id = 1; }
message ProjectListReturned { repeated string project_ids = 1; }
message ProjectCreateRequested { string display_name = 1; }
message ProjectCreated { string project_id = 1; string display_name = 2; }
message ProjectOpenRequested { string project_id = 1; }
message ProjectOpened { string project_id = 1; string current_requirements_revision = 2; }
message ProjectDeleteRequested { string project_id = 1; string confirmation = 2; }
message ProjectDeleted { string project_id = 1; }

message ProviderCredentialSubmitted {
  string provider = 1;
  bytes api_key = 2;
  string retention = 3;
}
message ProviderCredentialAccepted { string credential_ref = 1; string billing_source = 2; }
message ProviderCredentialRejected { string reason = 1; }

message ChatPromptSubmitted {
  string project_id = 1;
  string conversation_id = 2;
  string message_id = 3;
  repeated ChatContentPart parts = 4;
  repeated string attachment_ids = 5;
  string credential_ref = 6;
  string model = 7;
}
message ModelRequestIssued {
  string project_id = 1;
  string conversation_id = 2;
  string request_id = 3;
  string model = 4;
  string research_mode = 5;
}
message ModelResponseReturned {
  string request_id = 1;
  repeated ChatContentPart parts = 2;
  uint64 input_tokens = 3;
  uint64 cached_input_tokens = 4;
  uint64 output_tokens = 5;
  uint64 context_window_tokens = 6;
}
message ChatResponseRendered {
  string conversation_id = 1;
  string message_id = 2;
  repeated ChatContentPart parts = 3;
  repeated string artifact_ids = 4;
  uint64 input_tokens = 5;
  uint64 cached_input_tokens = 6;
  uint64 output_tokens = 7;
  uint64 context_window_tokens = 8;
  string usage_warning = 9;
  uint64 response_time_ms = 10;
}
message ChatTurnFailed { string request_id = 1; string reason = 2; }
message ConversationCheckpointRequested { string project_id = 1; string conversation_id = 2; string through_message_id = 3; }
message ConversationCheckpointBlocked { string project_id = 1; string conversation_id = 2; string checkpoint_id = 3; string reason_code = 4; repeated string blocking_artifact_ids = 5; }
message ConversationCheckpointCommitted { string project_id = 1; string conversation_id = 2; string through_message_id = 3; repeated string artifact_revisions = 4; string manifest_digest = 5; }
message ConversationHistoryCleared { string project_id = 1; string conversation_id = 2; string cleared_through_message_id = 3; uint64 retained_message_count = 4; }

message MermaidDraftSaved {
  string project_id = 1;
  string artifact_id = 2;
  uint64 revision = 3;
  string diagram_kind = 4;
  string source = 5;
}
message MermaidPreviewRendered { string artifact_id = 1; uint64 revision = 2; bytes svg = 3; }
message MediaUploaded {
  string project_id = 1;
  string artifact_id = 2;
  string media_type = 3;
  bytes content = 4;
}
message MediaRendered { string artifact_id = 1; string media_type = 2; string content_url = 3; }
message ArtifactPromoted { string artifact_id = 1; uint64 revision = 2; string target_status = 3; }
message CandidateRequirementUpdated { string requirements_revision = 1; repeated string source_artifact_ids = 2; }

message ArtifactListRequested {
  string project_id = 1;
  string kind_filter = 2;
  string scenario_filter = 3;
  string status_filter = 4;
}
message ArtifactSummary {
  string artifact_id = 1;
  uint64 current_revision = 2;
  uint64 accepted_revision = 3;
  string status = 4;
  string validation = 5;
  string media_type = 6;
  string title = 7;
  string scenario_id = 8;
  string diagram_kind = 9;
  string updated_by = 10;
  uint64 updated_at_ms = 11;
}
message ArtifactListReturned { repeated ArtifactSummary artifacts = 1; }
message ArtifactOpenRequested { string artifact_id = 1; uint64 revision = 2; }
message ArtifactOpened {
  string artifact_id = 1;
  uint64 revision = 2;
  string status = 3;
  string media_type = 4;
  bytes content = 5;
  string provenance = 6;
}
\end{lstlisting}


Envelope metadata remain distinct. \texttt{prior\_ids} stores immutable
predecessor links; any \texttt{map[Event][]Event} successor index is derived,
not serialized as mutable truth. \texttt{session\_id} separates concurrent
conversations, \texttt{task\_id} selects the grammar, and
\texttt{src}/\texttt{dst} must match the message schema. As defined above,
\texttt{cause\_id} is a separate acyclic direct-cause relation usable by roots,
intermediate messages, nested task boundaries, and terminals. Its tree groups
renderer boxes without ordering concurrent children.

\section{Complete LeanFM task grammar}

The following is the whole accepted ChatUI conversation grammar, rendered from
the typed \texttt{TaskGrammar} values. Pronounce
\texttt{A ! B : Message} as ``A sends Message to B.'' \texttt{then} is ordered
composition; \texttt{choose/or} is an explicit alternative;
\texttt{optional} permits zero or one occurrence; \texttt{fork/join all} waits
for every branch; and \texttt{race/join first} accepts the first completion and
cancels the remaining branches. A message not admitted by the selected task's
current residual grammar is illegal even if its protobuf bytes decode.

\section{Authenticate an individual user}
\noindent\textbf{Task:} \texttt{authenticate\_user}. Entry: \texttt{signed\_out}; terminals: \texttt{authenticated, rejected}.
\begin{verbatim}
UserBrowser ! LeanFMServer : LoginSubmitted
then
choose
  LeanFMServer ! UserBrowser : LoginAccepted
or
  LeanFMServer ! UserBrowser : LoginRejected
\end{verbatim}
\section{Create an account using a user-supplied provider API key}
\noindent\textbf{Task:} \texttt{create\_account}. Entry: \texttt{registration}; terminals: \texttt{created, rejected}.
\begin{verbatim}
UserBrowser ! LeanFMServer : AccountCreateSubmitted
then
LeanFMServer ! ModelProvider : ProviderKeyCheckRequested
then
choose
  ModelProvider ! LeanFMServer : ProviderKeyVerified
  then
  LeanFMServer ! UserBrowser : AccountCreated
or
  ModelProvider ! LeanFMServer : ProviderKeyRejected
  then
  LeanFMServer ! UserBrowser : AccountCreateRejected
\end{verbatim}
\section{List, create, open, and delete owned projects}
\noindent\textbf{Task:} \texttt{manage\_projects}. Entry: \texttt{ready}; terminals: \texttt{listed, created, opened, deleted}.
\begin{verbatim}
choose
  UserBrowser ! LeanFMServer : ProjectListRequested
  then
  LeanFMServer ! UserBrowser : ProjectListReturned
or
  UserBrowser ! LeanFMServer : ProjectCreateRequested
  then
  LeanFMServer ! UserBrowser : ProjectCreated
or
  UserBrowser ! LeanFMServer : ProjectOpenRequested
  then
  LeanFMServer ! UserBrowser : ProjectOpened
or
  UserBrowser ! LeanFMServer : ProjectDeleteRequested
  then
  LeanFMServer ! UserBrowser : ProjectDeleted
\end{verbatim}
\section{Turn a fresh project conversation into accepted requirements and implementation}
\noindent\textbf{Task:} \texttt{develop\_project\_specification}. Entry: \texttt{project\_loaded}; terminals: \texttt{ready\_to\_implement}.
\begin{verbatim}
LeanFMServer ! UserBrowser : ProblemInquiryPresented
then
UserBrowser ! LeanFMServer : ProblemStatementSubmitted
then
repeat
  choose
    empty
  or
    LeanFMServer ! UserBrowser : ProjectBookRevised
  then
  LeanFMServer ! UserBrowser : AmbiguityQuestionPresented
  then
  UserBrowser ! LeanFMServer : ClarificationSubmitted
then
choose
  empty
or
  LeanFMServer ! UserBrowser : ProjectBookRevised
then
LeanFMServer ! UserBrowser : RequirementsReviewPresented
then
UserBrowser ! LeanFMServer : SpecificationRequirementsAccepted
then
LeanFMServer ! UserBrowser : ImplementationReviewPresented
then
UserBrowser ! LeanFMServer : SpecificationImplementationAccepted
then
LeanFMServer ! UserBrowser : ImplementationAuthorizationConfirmed
\end{verbatim}
\section{Select user-funded or platform-funded model access}
\noindent\textbf{Task:} \texttt{configure\_provider}. Entry: \texttt{unconfigured}; terminals: \texttt{configured, rejected}.
\begin{verbatim}
UserBrowser ! LeanFMServer : ProviderCredentialSubmitted
then
choose
  LeanFMServer ! UserBrowser : ProviderCredentialAccepted
or
  LeanFMServer ! UserBrowser : ProviderCredentialRejected
\end{verbatim}
\section{Conduct one durable project-scoped requirements turn}
\noindent\textbf{Task:} \texttt{argue\_requirements}. Entry: \texttt{ready}; terminals: \texttt{rendered, failed}.
\begin{verbatim}
UserBrowser ! LeanFMServer : ChatPromptSubmitted
then
LeanFMServer ! ModelProvider : ModelRequestIssued
then
choose
  ModelProvider ! LeanFMServer : ModelResponseReturned
  then
  LeanFMServer ! UserBrowser : ChatResponseRendered
or
  LeanFMServer ! UserBrowser : ChatTurnFailed
\end{verbatim}
\section{Compact a project conversation into durable specifications before clearing it}
\noindent\textbf{Task:} \texttt{checkpoint\_conversation}. Entry: \texttt{retained}; terminals: \texttt{cleared, blocked}.
\begin{verbatim}
UserBrowser ! LeanFMServer : ConversationCheckpointRequested
then
choose
  LeanFMServer ! UserBrowser : ConversationCheckpointBlocked
or
  LeanFMServer ! UserBrowser : ConversationCheckpointCommitted
  then
  LeanFMServer ! UserBrowser : ConversationHistoryCleared
\end{verbatim}
\section{Sketch and promote multimodal requirement artifacts}
\noindent\textbf{Task:} \texttt{markerboard}. Entry: \texttt{editing}; terminals: \texttt{candidate}.
\begin{verbatim}
choose
  UserBrowser ! LeanFMServer : MermaidDraftSaved
  then
  LeanFMServer ! UserBrowser : MermaidPreviewRendered
  then
  UserBrowser ! LeanFMServer : ArtifactPromoted
  then
  LeanFMServer ! UserBrowser : CandidateRequirementUpdated
or
  UserBrowser ! LeanFMServer : MediaUploaded
  then
  LeanFMServer ! UserBrowser : MediaRendered
  then
  UserBrowser ! LeanFMServer : ArtifactPromoted
  then
  LeanFMServer ! UserBrowser : CandidateRequirementUpdated
\end{verbatim}
\section{Inspect current, accepted, superseded, and stale artifacts}
\noindent\textbf{Task:} \texttt{inspect\_artifacts}. Entry: \texttt{closed}; terminals: \texttt{listed, opened}.
\begin{verbatim}
UserBrowser ! LeanFMServer : ArtifactListRequested
then
LeanFMServer ! UserBrowser : ArtifactListReturned
then
choose
  empty
or
  UserBrowser ! LeanFMServer : ArtifactOpenRequested
  then
  LeanFMServer ! UserBrowser : ArtifactOpened
\end{verbatim}
\section{Inspect one task's linked interaction, state machine, constraints, and protobuf events}
\noindent\textbf{Task:} \texttt{inspect\_task\_model}. Entry: \texttt{closed}; terminals: \texttt{linked, sequence\_visible, object\_visible}.
\begin{verbatim}
UserBrowser ! LeanFMServer : TaskModelRequested
then
LeanFMServer ! UserBrowser : TaskModelReturned
then
choose
  empty
or
  UserBrowser ! LeanFMServer : MessageSequenceRequested
  then
  LeanFMServer ! UserBrowser : MessageSequenceReturned
  then
  choose
    empty
  or
    UserBrowser ! LeanFMServer : ProtobufObjectRequested
    then
    LeanFMServer ! UserBrowser : ProtobufObjectReturned
\end{verbatim}
\section{Search the public web, follow supplied HTTP URLs, and return clickable citations}
\noindent\textbf{Task:} \texttt{research\_web}. Entry: \texttt{ready}; terminals: \texttt{returned, failed}.
\begin{verbatim}
UserBrowser ! LeanFMServer : WebResearchRequested
then
LeanFMServer ! ModelProvider : WebResearchDispatched
then
choose
  ModelProvider ! LeanFMServer : WebSearchPerformed
  then
  ModelProvider ! LeanFMServer : WebPageOpened
  then
  LeanFMServer ! UserBrowser : WebResearchReturned
or
  LeanFMServer ! UserBrowser : WebResearchFailed
\end{verbatim}
\section{Recall diagram syntax or propose AI-generated detail for one markerboard revision}
\noindent\textbf{Task:} \texttt{assist\_markerboard}. Entry: \texttt{editing}; terminals: \texttt{helped, updated, rejected}.
\begin{verbatim}
UserBrowser ! LeanFMServer : MarkerboardAssistRequested
then
choose
  LeanFMServer ! UserBrowser : MarkerboardSyntaxHelpReturned
or
  LeanFMServer ! ModelProvider : MarkerboardAssistDispatched
  then
  ModelProvider ! LeanFMServer : MarkerboardAssistReturned
  then
  LeanFMServer ! UserBrowser : MarkerboardPatchProposed
  then
  choose
    UserBrowser ! LeanFMServer : MarkerboardPatchAccepted
    then
    LeanFMServer ! UserBrowser : MarkerboardDraftUpdated
  or
    UserBrowser ! LeanFMServer : MarkerboardPatchRejected
\end{verbatim}


\section{The grammar as an executable canvas}

The complete grammar is one navigable canvas rather than one permanently huge
drawing. A task or production node may be closed to its typed boundary or
opened to show its interaction messages, residual states, parameters,
constraints, paired measurements and nested productions. Opening and closing
is presentation state only. It never changes the grammar or conceals a failing
obligation.

Each editable parameter has one typed definition and explicit dependency
edges. Changing a queue capacity, actor population, timeout, probability or
other parameter marks every transitive dependent stale and regenerates the
affected productions, reducers, proofs, codecs, charts and views. The canvas
shows that cascade before the revised artifact can be accepted.

A production known to cross HTTP may attach its transport renderer at that
production. The binding names the method, path template, request and response
protobuf types, status alternatives, content type, byte encoder and byte
decoder. Generated fixtures emit the exact request and response bytes and read
them back. Conformance requires decoded emitted values to equal their originals
and their message atoms to be accepted by the selected residual grammar. Thus
the high-level canvas can be progressively opened all the way down to bytes
without confusing transport with the message value.

\section{What is checked, and what remains}

Lean elaboration checks the actors, message schemas, task states, transitions,
grammars, constraints, desired-output provenance, and implementation mappings.
The chapter generator rejects an accepted HTTP request with no named response
mapping and rejects incomplete running-route records. The ordinary repository
checks also validate the protobuf source against the typed requirements.

This does not yet prove that every accepted route has a handler: the running
table makes that gap explicit. Nor does the current server yet append every
browser exchange to the specified binary event journal. Completing those
handlers and the schema-specific codec remains implementation work. Until then,
the accepted grammar is the protocol contract, the running table is the honest
conformance statement, and neither may be substituted for the other.
